% !TeX program = xelatex
% Компиляция: xelatex Course_Work_Report.tex (два запуска)
% Альтернатива: tectonic Course_Work_Report.tex
\documentclass[14pt,a4paper]{extarticle}
\usepackage{fontspec}
\setmainfont{Times New Roman}
\setmonofont{DejaVuSansMono.ttf}[BoldFont=DejaVuSansMono-Bold.ttf,Scale=0.82]
\usepackage[english]{babel}
\babelprovide[import,main]{russian}
\usepackage{amsmath,unicode-math}
\setmathfont{texgyretermes-math.otf}
\usepackage[left=30mm,right=15mm,top=20mm,bottom=20mm]{geometry}
\usepackage{setspace,indentfirst}
\onehalfspacing
\setlength{\parindent}{1.25cm}
\setlength{\parskip}{0pt}
\usepackage{titlesec}
\titleformat{\section}{\normalfont\bfseries}{\thesection}{0.6em}{}
\titleformat{\subsection}{\normalfont\bfseries}{\thesubsection}{0.6em}{}
\titlespacing*{\section}{\parindent}{0pt}{12pt}
\titlespacing*{\subsection}{\parindent}{12pt}{6pt}
\newcommand{\sectionbreak}{\clearpage}
\usepackage{graphicx,float,caption,array,tabularx}
\DeclareCaptionLabelSeparator{dash}{~---~}
\captionsetup{labelsep=dash,font=normalsize,skip=6pt}
\captionsetup[figure]{name=Рисунок,justification=centering}
\captionsetup[table]{name=Таблица,justification=raggedright,singlelinecheck=false,position=top}
\renewcommand{\arraystretch}{1.15}
\newcolumntype{L}{>{\raggedright\arraybackslash}X}
\usepackage{fvextra,newfloat,needspace}
\DeclareFloatingEnvironment[name=Листинг]{listing}
\captionsetup[listing]{hypcap=false,justification=raggedright,singlelinecheck=false}
\fvset{numbers=left,numbersep=5pt,fontsize=\scriptsize,frame=single,
  breaklines=true,breakanywhere=true,breaksymbolleft={},tabsize=4}
\usepackage[hidelinks]{hyperref}
\newcommand{\code}[1]{\texttt{\detokenize{#1}}}
\newsavebox{\functiondescription}
\newcommand{\fn}[3]{\par
  \setbox\functiondescription=\vbox{\hsize=\linewidth
  \noindent\textbf{\code{#1}}. \textit{Вход:} #2\par
  \noindent\textit{Выход:} #3\par}
  \Needspace{\dimexpr\ht\functiondescription+\dp\functiondescription+6\baselineskip\relax}
  \medskip\unvbox\functiondescription}
\newcommand{\structsection}[1]{\clearpage\phantomsection
  \addcontentsline{toc}{section}{#1}
  \begin{center}\textbf{#1}\end{center}\vspace{6pt}}
\pagestyle{plain}
\sloppy

\begin{document}
\begin{titlepage}
\begin{singlespace}
\begin{center}
Министерство науки и высшего образования Российской Федерации\\
Санкт-Петербургский политехнический университет Петра Великого\\
Институт компьютерных наук и кибербезопасности\\
Высшая школа технологий искусственного интеллекта
\end{center}
\vspace{5cm}
\begin{center}
\textbf{ОТЧЁТ}\\
по курсовой работе\\[6pt]
«Конечная арифметика. Вариант X8»\\[6pt]
по дисциплине «Дискретная математика»
\end{center}
\vspace{3.5cm}
\begin{flushright}
\begin{tabular}{@{}l@{}}
Выполнил\\
студент гр. 5130201/50302 \underline{\hspace{2.5cm}} Р. Ю. Дужнов\\[12pt]
Проверил\\
преподаватель \underline{\hspace{2.5cm}} А. В. Востров\\[12pt]
« \underline{\hspace{1cm}} » \underline{\hspace{3cm}} 2026 г.
\end{tabular}
\end{flushright}
\vfill
\begin{center}Санкт-Петербург\\2026\end{center}
\end{singlespace}
\end{titlepage}
\setcounter{page}{2}
\tableofcontents

\structsection{ВВЕДЕНИЕ}

В курсовой работе требуется разработать программу для выполнения арифметических операций над числами, записанными с помощью символов конечного алфавита. Для варианта X8 заданы элементы $a,b,c,d,e,f,g,h$ и правило перехода к следующему элементу. На основе этого правила необходимо определить порядок цифр, составить таблицы операций и организовать вычисления над многоразрядными числами.

Программа должна выполнять сложение, вычитание, умножение и целочисленное деление с остатком. При вычислениях требуется учитывать знак числа, переносы и заимствования между разрядами. Для исходных чисел используется ограничение в восемь разрядов; при превышении этого ограничения результатом программа должна сообщать о переполнении разрядной сетки.

Цель работы --- реализовать позиционную арифметику, в которой действия над отдельными цифрами выполняются по таблицам, без перевода операндов в стандартный целочисленный тип. Для достижения цели необходимо определить математическую модель, разработать алгоритмы операций, реализовать консольное меню и проверить правильность вычислений.

Пользователь должен иметь возможность вводить выражения, просматривать таблицы операций и повторять вычисления без перезапуска программы. Требуется предусмотреть обработку недопустимых символов, ошибочного формата выражения, превышения длины операнда и деления на ноль.

\section{Математическое описание}

\subsection{Алфавит и порядок цифр}

Исходный алфавит содержит восемь символов:
\[
X=\{a,b,c,d,e,f,g,h\}.
\]
Правило перехода «$+1$» имеет вид
\[
\begin{gathered}
a\to b,\quad b\to e,\quad c\to a,\quad d\to a,\\
e\to h,\quad f\to g,\quad g\to\{c,d\},\quad h\to f.
\end{gathered}
\]

В принятой в программе модели символы $c$ и $d$ обозначают одну цифру: $c\sim d$. Для внутреннего представления выбирается $c$; все введённые символы $d$ заменяются на $c$. Цепочка последовательных цифр принимает вид
\[
a\to b\to e\to h\to f\to g\to c\to a.
\]
Таким образом, система содержит семь различных значений цифр. Нулевая цифра --- $a$, единичная --- $b$, основание позиционной системы счисления равно семи. Порядок цифр:
\[
a<b<e<h<f<g<c.
\]

Для математической проверки удобно ввести функцию значения цифры $\nu$ (таблица~\ref{tab:digits}). В самой программе перевода многоразрядных операндов в десятичную систему нет.

\begin{table}[H]
\caption{Соответствие цифр их значениям}\label{tab:digits}
\centering
\begin{tabular}{|c|c|c|c|c|c|c|c|}\hline
Символ & $a$ & $b$ & $e$ & $h$ & $f$ & $g$ & $c$, $d$\\\hline
$\nu$ & 0 & 1 & 2 & 3 & 4 & 5 & 6\\\hline
\end{tabular}
\end{table}

Важно различать циклический переход между цифрами и сложение позиционных чисел. После максимальной цифры $c$ младшая цифра становится $a$, а в следующий разряд переносится $b$. Поэтому для чисел $c+b=ba$, тогда как значение табличной функции следующей цифры равно $a$.

\subsection{Позиционная запись числа}

Пусть $x_{k-1}\ldots x_1x_0$ --- запись модуля числа, где $x_0$ является младшей цифрой. Значение записи определяется формулой
\[
V(x_{k-1}\ldots x_0)=\sum_{i=0}^{k-1}\nu(x_i)7^i.
\]
Знак хранится отдельно от модуля. Например,
\[
V(bc)=1\cdot7+6=13,\qquad V(ea)=2\cdot7+0=14.
\]
Ведущие нулевые цифры не изменяют значение числа и удаляются. Запись, состоящая только из символов $a$, приводится к одному символу $a$. У нулевого результата отрицательный знак снимается.

Сначала модули чисел сравниваются по длине нормализованных записей. Если длины различаются, больше модуль с более длинной записью. При одинаковой длине результат определяет старший разряд, в котором цифры различаются, с учётом порядка $a<b<e<h<f<g<c$.

Входное число содержит от одного до восьми буквенных разрядов. Максимальный допустимый модуль равен
\[
V(cccccccc)=7^8-1=5\,764\,800.
\]
Ограничение проверяется до удаления ведущих нулей. Поэтому строка из девяти символов $a$ также отклоняется. Для результата длина не ограничивается восемью разрядами: программа сохраняет полную запись и выводит предупреждение.

\subsection{Таблицы одноразрядных операций}

Пусть $x,y$ --- цифры и $\delta=\nu^{-1}$ --- обратное отображение в канонический алфавит. Таблицы младших цифр задаются следующим образом:
\[
\begin{aligned}
\operatorname{ADD}(x,y)&=\delta((\nu(x)+\nu(y))\bmod7),\\
\operatorname{SUB}(x,y)&=\delta((\nu(x)-\nu(y))\bmod7),\\
\operatorname{MUL}(x,y)&=\delta((\nu(x)\nu(y))\bmod7).
\end{aligned}
\]
При вычитании берётся неотрицательный остаток по модулю семи. Старшая цифра произведения определяется отдельно:
\[
\operatorname{MULHI}(x,y)=\delta\!\left(\left\lfloor\frac{\nu(x)\nu(y)}7\right\rfloor\right).
\]
Для деления на ненулевую цифру таблица хранит пару «частное, остаток»:
\[
\operatorname{DIV}(x,y)=
\left(\delta\!\left(\left\lfloor\frac{\nu(x)}{\nu(y)}\right\rfloor\right),
\delta(\nu(x)\bmod\nu(y))\right).
\]
Таблицы содержат только канонические символы. Строка соответствует первому операнду, столбец --- второму. В отчёте цифры расположены по возрастанию значения; в консоли программа выводит их в порядке ключей \code{std::map}, то есть $a,b,c,e,f,g,h$.

\begin{table}[H]
\caption{Сложение: младшая цифра результата}\label{tab:add}
\centering
\begin{tabular}{|c|c|c|c|c|c|c|c|}\hline
$+$ & $a$ & $b$ & $e$ & $h$ & $f$ & $g$ & $c$\\\hline
$a$ & a & b & e & h & f & g & c\\\hline
$b$ & b & e & h & f & g & c & a\\\hline
$e$ & e & h & f & g & c & a & b\\\hline
$h$ & h & f & g & c & a & b & e\\\hline
$f$ & f & g & c & a & b & e & h\\\hline
$g$ & g & c & a & b & e & h & f\\\hline
$c$ & c & a & b & e & h & f & g\\\hline
\end{tabular}
\end{table}

\begin{table}[H]
\caption{Вычитание: разность по модулю семи}\label{tab:sub}
\centering
\begin{tabular}{|c|c|c|c|c|c|c|c|}\hline
$-$ & $a$ & $b$ & $e$ & $h$ & $f$ & $g$ & $c$\\\hline
$a$ & a & c & g & f & h & e & b\\\hline
$b$ & b & a & c & g & f & h & e\\\hline
$e$ & e & b & a & c & g & f & h\\\hline
$h$ & h & e & b & a & c & g & f\\\hline
$f$ & f & h & e & b & a & c & g\\\hline
$g$ & g & f & h & e & b & a & c\\\hline
$c$ & c & g & f & h & e & b & a\\\hline
\end{tabular}
\end{table}

\begin{table}[H]
\caption{Умножение: младшая цифра произведения}\label{tab:mul}
\centering
\begin{tabular}{|c|c|c|c|c|c|c|c|}\hline
$\cdot$ & $a$ & $b$ & $e$ & $h$ & $f$ & $g$ & $c$\\\hline
$a$ & a & a & a & a & a & a & a\\\hline
$b$ & a & b & e & h & f & g & c\\\hline
$e$ & a & e & f & c & b & h & g\\\hline
$h$ & a & h & c & e & g & b & f\\\hline
$f$ & a & f & b & g & e & c & h\\\hline
$g$ & a & g & h & b & c & f & e\\\hline
$c$ & a & c & g & f & h & e & b\\\hline
\end{tabular}
\end{table}

\begin{table}[H]
\caption{Умножение: старшая цифра произведения}\label{tab:mulhi}
\centering
\begin{tabular}{|c|c|c|c|c|c|c|c|}\hline
$\cdot$ & $a$ & $b$ & $e$ & $h$ & $f$ & $g$ & $c$\\\hline
$a$ & a & a & a & a & a & a & a\\\hline
$b$ & a & a & a & a & a & a & a\\\hline
$e$ & a & a & a & a & b & b & b\\\hline
$h$ & a & a & a & b & b & e & e\\\hline
$f$ & a & a & b & b & e & e & h\\\hline
$g$ & a & a & b & e & e & h & f\\\hline
$c$ & a & a & b & e & h & f & g\\\hline
\end{tabular}
\end{table}

\begin{table}[H]
\caption{Деление цифр: частное и остаток}\label{tab:div}
\centering
\begin{tabular}{|c|c|c|c|c|c|c|}\hline
$/ $ & $b$ & $e$ & $h$ & $f$ & $g$ & $c$\\\hline
$a$ & $(a,a)$ & $(a,a)$ & $(a,a)$ & $(a,a)$ & $(a,a)$ & $(a,a)$\\\hline
$b$ & $(b,a)$ & $(a,b)$ & $(a,b)$ & $(a,b)$ & $(a,b)$ & $(a,b)$\\\hline
$e$ & $(e,a)$ & $(b,a)$ & $(a,e)$ & $(a,e)$ & $(a,e)$ & $(a,e)$\\\hline
$h$ & $(h,a)$ & $(b,b)$ & $(b,a)$ & $(a,h)$ & $(a,h)$ & $(a,h)$\\\hline
$f$ & $(f,a)$ & $(e,a)$ & $(b,b)$ & $(b,a)$ & $(a,f)$ & $(a,f)$\\\hline
$g$ & $(g,a)$ & $(e,b)$ & $(b,e)$ & $(b,b)$ & $(b,a)$ & $(a,g)$\\\hline
$c$ & $(c,a)$ & $(h,a)$ & $(e,a)$ & $(b,e)$ & $(b,b)$ & $(b,a)$\\\hline
\end{tabular}
\end{table}

Например, ячейка $\operatorname{ADD}(c,b)=a$ задаёт только младшую цифру суммы; перенос равен $b$. Для произведения $c\cdot c$ таблицы дают младшую цифру $b$ и старшую $g$, то есть результат $gb$. Ячейка деления $(c,e)$ содержит $(h,a)$: $6=3\cdot2+0$.

\subsection{Многоразрядные операции и знаки}

Сложение модулей выполняется справа налево. Сначала складываются две цифры, затем к полученной цифре прибавляется перенос предыдущего разряда. Перенос возникает, если после табличного сложения цифра уменьшается относительно первого слагаемого в установленном порядке. Оба возможных переноса объединяются логическим «или».

Вычитание модулей выполняется при условии $A\geq B$. Для каждого разряда вычисляется табличная разность цифр и вычитается заимствование. Если цифры уменьшаемого недостаточно, из следующего разряда заимствуется единица. После обработки всех разрядов ведущие нули удаляются.

Умножение организовано «в столбик». Для каждой цифры второго множителя вычисляется произведение первого множителя на эту цифру. Таблицы \code{MUL} и \code{MULHI} дают младшую цифру и перенос. Умножение на $7^i$ выполняется дописыванием $i$ нулевых цифр $a$ справа. Частичные произведения суммируются.

Деление модулей выполняется слева направо. На каждом шаге к текущему остатку дописывается очередная цифра делимого. Пока остаток не меньше делителя, делитель вычитается, а цифра частного увеличивается на единицу. До приписывания следующей цифры остаток меньше делителя; следовательно, новая цифра частного не превышает шести. Для двух одноразрядных модулей используется таблица \code{DIV}.

При сложении чисел одинакового знака складываются их модули и сохраняется общий знак. При разных знаках из большего модуля вычитается меньший, а результат получает знак большего по модулю операнда. Вычитание сводится к сложению с противоположным числом. Знак произведения отрицателен, если знаки множителей различаются.

Для деления со знаком используется евклидово условие:
\[
n=q\,m+r,\qquad 0\leq r<|m|,\quad m\neq0.
\]
Сначала находятся $Q,R$ из равенства $|n|=Q|m|+R$. Если $n<0$ и $R\neq0$, модуль частного увеличивается на единицу, а остаток заменяется на $|m|-R$. Знак частного определяется различием знаков операндов. Например,
\[
-13=(-7)\cdot2+1=7\cdot(-2)+1.
\]
В буквенной записи это соответствует частному $-ba$ при делении $-bc$ на $e$ и частному $ba$ при делении $-bc$ на $-e$; остаток в обоих случаях равен $b$.

\section{Особенности реализации}

Программа реализована на языке C++ в файле \code{course-work/main.cpp}. Таблицы \code{ADD}, \code{SUB}, \code{MUL} и \code{MULHI} имеют тип \code{map<char, map<char, char>>}. Таблица \code{DIV} хранит для каждой пары цифр пару символов частного и остатка. Эти таблицы заданы явно и не строятся через перевод многоразрядных чисел в десятичную систему.

Модуль числа представлен строкой \code{Num}, знак и модуль --- структурой \code{SNum}. Тип \code{QR} объединяет частное и остаток деления модулей. Перечисление \code{Cmp} содержит результаты сравнения \code{LT}, \code{EQ}, \code{GT}. Константы \code{ZERO}, \code{ONE} и \code{MAX_DIGITS} равны соответственно \code{'a'}, \code{'b'} и восьми.

\begin{singlespace}
\captionof{listing}{Типы данных и основные константы}
\begin{Verbatim}
using Num = string;
struct SNum { bool neg; Num mag; };
enum class Cmp { LT, EQ, GT };
struct QR { Num q, r; };
static const char ZERO = 'a';
static const char ONE = 'b';
static const size_t MAX_DIGITS = 8;
\end{Verbatim}
\end{singlespace}

\subsection{Порядок цифр и нормализация модулей}

\par\noindent\begin{minipage}{\linewidth}
\fn{succ}{\code{x} --- каноническая цифра.}{Следующая цифра циклического порядка. Используется ячейка сложения с \code{ONE}; после $c$ возвращается $a$.}
\begin{singlespace}
\captionof{listing}{\code{succ}}
\begin{Verbatim}
static char succ(char x) { return ADD.at(x).at(ONE); }
\end{Verbatim}
\end{singlespace}
\end{minipage}\par\medskip

\par\noindent\begin{minipage}{\linewidth}
\fn{buildOrder}{Таблица \code{ADD} и константы \code{ZERO}, \code{ONE}.}{Заполнено отношение \code{LESS}: для каждой цифры отмечены все цифры, стоящие после неё в цикле от $a$. Например, устанавливаются $b<e$ и $g<c$. Вызывается один раз в начале работы.}
\begin{singlespace}
\captionof{listing}{\code{buildOrder}}
\begin{Verbatim}
static void buildOrder() {
    string seen;                       // уже пройденные элементы цепочки
    char cur = ZERO;
    do {
        for (char s : seen) LESS[s][cur] = true;   // всё, что раньше cur, меньше cur
        seen.push_back(cur);
        cur = succ(cur);
    } while (cur != ZERO);
}
\end{Verbatim}
\end{singlespace}
\end{minipage}\par\medskip

\par\noindent\begin{minipage}{\linewidth}
\fn{lessD}{\code{x}, \code{y} --- канонические цифры.}{\code{true}, если $x<y$ в установленном порядке, иначе \code{false}. Сравнение кодов символов не используется.}
\begin{singlespace}
\captionof{listing}{\code{lessD}}
\begin{Verbatim}
static bool lessD(char x, char y) {
    auto it = LESS.find(x);
    if (it == LESS.end()) return false;
    auto jt = it->second.find(y);
    return jt != it->second.end() && jt->second;
}
\end{Verbatim}
\end{singlespace}
\end{minipage}\par\medskip

\par\noindent\begin{minipage}{\linewidth}
\fn{strip}{\code{s} --- строковая запись модуля.}{Нормализованный модуль без ведущих символов $a$. Для нуля возвращается строка \code{"a"}. Например, \code{"aabc"} превращается в \code{"bc"}.}
\begin{singlespace}
\captionof{listing}{\code{strip}}
\begin{Verbatim}
static Num strip(Num s) {
    auto it = s.begin();
    while (it != s.end() && *it == ZERO) ++it;
    s.erase(s.begin(), it);
    if (s.empty()) s = string(1, ZERO);
    return s;
}
\end{Verbatim}
\end{singlespace}
\end{minipage}\par\medskip

\par\noindent\begin{minipage}{\linewidth}
\fn{cmpMag}{\code{A}, \code{B} --- нормализованные модули.}{Результат сравнения. Разряды просматриваются от младших к старшим; каждое более старшее различие заменяет предыдущий результат. Если длины различаются, больше более длинная запись.}
\begin{singlespace}
\captionof{listing}{\code{cmpMag}}
\begin{Verbatim}
static Cmp cmpMag(const Num& A, const Num& B) {
    auto ia = A.rbegin();
    auto ib = B.rbegin();
    Cmp res = Cmp::EQ;
    while (ia != A.rend() && ib != B.rend()) {
        if (*ia != *ib) res = lessD(*ia, *ib) ? Cmp::LT : Cmp::GT;
        ++ia; ++ib;
    }
    if (ia != A.rend()) return Cmp::GT;
    if (ib != B.rend()) return Cmp::LT;
    return res;
}
\end{Verbatim}
\end{singlespace}
\end{minipage}\par\medskip

\Needspace{20\baselineskip}
\subsection{Операции над модулями}

\par\noindent\begin{minipage}{\linewidth}
\fn{addMag}{\code{A}, \code{B} --- нормализованные модули.}{Нормализованная сумма модулей. Переносы \code{c1}, \code{c2} обнаруживают переход через максимальную цифру при двух табличных сложениях. Например, $c+b=ba$.}
\begin{singlespace}
\captionof{listing}{\code{addMag}}
\begin{Verbatim}
static Num addMag(const Num& A, const Num& B) {
    Num res;
    char carry = ZERO;
    auto ia = A.rbegin();
    auto ib = B.rbegin();
    while (ia != A.rend() || ib != B.rend() || carry != ZERO) {
        char x = (ia != A.rend()) ? *ia : ZERO;
        char y = (ib != B.rend()) ? *ib : ZERO;
        char s1 = ADD.at(x).at(y);
        bool c1 = lessD(s1, x);
        char s2 = ADD.at(s1).at(carry);
        bool c2 = lessD(s2, s1);
        res.insert(res.begin(), s2);
        carry = (c1 || c2) ? ONE : ZERO;
        if (ia != A.rend()) ++ia;
        if (ib != B.rend()) ++ib;
    }
    return strip(res);
}
\end{Verbatim}
\end{singlespace}
\end{minipage}\par\medskip

\par\noindent\begin{minipage}{\linewidth}
\fn{subMag}{\code{A}, \code{B} --- нормализованные модули, $A\geq B$.}{Нормализованная разность $A-B$. Переменная \code{borrow} содержит заимствование для следующего разряда. Предусловие обеспечивается вызывающими функциями.}
\begin{singlespace}
\captionof{listing}{\code{subMag}}
\begin{Verbatim}
static Num subMag(const Num& A, const Num& B) {
    Num res;
    char borrow = ZERO;
    auto ia = A.rbegin();
    auto ib = B.rbegin();
    while (ia != A.rend()) {
        char x = *ia;
        char y = (ib != B.rend()) ? *ib : ZERO;
        char d1 = SUB.at(x).at(y);
        bool b1 = lessD(x, y);
        char d2 = SUB.at(d1).at(borrow);
        bool b2 = lessD(d1, borrow);
        res.insert(res.begin(), d2);
        borrow = (b1 || b2) ? ONE : ZERO;
        ++ia;
        if (ib != B.rend()) ++ib;
    }
    return strip(res);
}
\end{Verbatim}
\end{singlespace}
\end{minipage}\par\medskip

\par\noindent\begin{minipage}{\linewidth}
\fn{mulDigit}{\code{A} --- модуль; \code{y} --- одна цифра.}{Произведение модуля на цифру. К младшей цифре произведения прибавляется предыдущий перенос; при переполнении старшая цифра увеличивается функцией \code{succ}.}
\begin{singlespace}
\captionof{listing}{\code{mulDigit}}
\begin{Verbatim}
static Num mulDigit(const Num& A, char y) {
    Num res;
    char carry = ZERO;
    for (auto ia = A.rbegin(); ia != A.rend(); ++ia) {
        char lo = MUL.at(*ia).at(y);
        char hi = MULHI.at(*ia).at(y);
        char t  = ADD.at(lo).at(carry);
        if (lessD(t, lo)) hi = succ(hi);       // сложение с переносом дало ещё одну единицу
        res.insert(res.begin(), t);
        carry = hi;
    }
    if (carry != ZERO) res.insert(res.begin(), carry);
    return strip(res);
}
\end{Verbatim}
\end{singlespace}
\end{minipage}\par\medskip

\par\noindent\begin{minipage}{\linewidth}
\fn{mulMag}{\code{A}, \code{B} --- нормализованные модули.}{Произведение модулей. Строка \code{shift} задаёт разрядный сдвиг: к частичному произведению дописываются символы $a$, после чего оно прибавляется к \code{total}.}
\begin{singlespace}
\captionof{listing}{\code{mulMag}}
\begin{Verbatim}
static Num mulMag(const Num& A, const Num& B) {
    Num total(1, ZERO);
    std::string shift;                         // сдвиг на разряд = дописать 'a' справа
    for (auto ib = B.rbegin(); ib != B.rend(); ++ib) {
        Num part = mulDigit(A, *ib);
        if (part != std::string(1, ZERO)) part += shift;
        total = addMag(total, part);
        shift.push_back(ZERO);
    }
    return total;
}
\end{Verbatim}
\end{singlespace}
\end{minipage}\par\medskip

\par\noindent\begin{minipage}{\linewidth}
\fn{divMag}{\code{N} --- модуль делимого; \code{D} --- ненулевой модуль делителя.}{Структура \code{QR} с частным и остатком, удовлетворяющими $N=qD+r$, $0\leq r<D$. Для однозначных операндов используется таблица; для остальных --- последовательное деление с повторным вычитанием.}
\begin{singlespace}
\captionof{listing}{\code{divMag}}
\begin{Verbatim}
static QR divMag(const Num& N, const Num& D) {
    if (N.size() == 1 && D.size() == 1) {      // обе - цифры: берём из таблицы DIV
        auto qr = DIV.at(N[0]).at(D[0]);
        return { std::string(1, qr.first), std::string(1, qr.second) };
    }
    Num q;
    Num R(1, ZERO);
    for (char d : N) {
        if (R == std::string(1, ZERO)) R = std::string(1, d);   // R = R*база + d : дописываем цифру справа
        else R.push_back(d);
        char cnt = ZERO;
        while (cmpMag(R, D) != Cmp::LT) {
            R = subMag(R, D);
            cnt = succ(cnt);
        }
        q.push_back(cnt);
    }
    return { strip(q), R };
}
\end{Verbatim}
\end{singlespace}
\end{minipage}\par\medskip

\subsection{Операции над числами со знаком}

\par\noindent\begin{minipage}{\linewidth}
\fn{isZero}{\code{m} --- нормализованный модуль.}{Признак равенства модуля нулю. Нормализованная запись нуля содержит ровно один символ $a$.}
\begin{singlespace}
\captionof{listing}{\code{isZero}}
\begin{Verbatim}
static bool isZero(const Num& m) { return m == std::string(1, ZERO); }
\end{Verbatim}
\end{singlespace}
\end{minipage}\par\medskip

\par\noindent\begin{minipage}{\linewidth}
\fn{fixZero}{\code{v} --- число со знаком.}{Число с тем же модулем; при нулевом модуле знак становится неотрицательным. Запись «минус ноль» в результатах исключается.}
\begin{singlespace}
\captionof{listing}{\code{fixZero}}
\begin{Verbatim}
static SNum fixZero(SNum v) { if (isZero(v.mag)) v.neg = false; return v; }
\end{Verbatim}
\end{singlespace}
\end{minipage}\par\medskip

\par\noindent\begin{minipage}{\linewidth}
\fn{addS}{\code{x}, \code{y} --- числа со знаком.}{Сумма чисел. Для одинаковых знаков вызывается \code{addMag}; для разных --- сравнение модулей и \code{subMag}. Равные модули с противоположными знаками дают ноль.}
\begin{singlespace}
\captionof{listing}{\code{addS}}
\begin{Verbatim}
static SNum addS(const SNum& x, const SNum& y) {
    SNum r;
    if (x.neg == y.neg) {
        r = { x.neg, addMag(x.mag, y.mag) };
    } else {
        Cmp c = cmpMag(x.mag, y.mag);
        if (c == Cmp::EQ)      r = { false, std::string(1, ZERO) };
        else if (c == Cmp::GT) r = { x.neg, subMag(x.mag, y.mag) };
        else                   r = { y.neg, subMag(y.mag, x.mag) };
    }
    return fixZero(r);
}
\end{Verbatim}
\end{singlespace}
\end{minipage}\par\medskip

\par\noindent\begin{minipage}{\linewidth}
\fn{subS}{\code{x}, \code{y} --- числа со знаком.}{Разность $x-y$, вычисленная как сумма $x+(-y)$. Нормализацию нулевого результата выполняет \code{addS}.}
\begin{singlespace}
\captionof{listing}{\code{subS}}
\begin{Verbatim}
static SNum subS(const SNum& x, const SNum& y) {
    return addS(x, SNum{ !y.neg, y.mag });
}
\end{Verbatim}
\end{singlespace}
\end{minipage}\par\medskip

\par\noindent\begin{minipage}{\linewidth}
\fn{mulS}{\code{x}, \code{y} --- числа со знаком.}{Произведение модулей со знаком, определяемым выражением \code{x.neg != y.neg}. У нулевого произведения знак снимается.}
\begin{singlespace}
\captionof{listing}{\code{mulS}}
\begin{Verbatim}
static SNum mulS(const SNum& x, const SNum& y) {
    return fixZero(SNum{ x.neg != y.neg, mulMag(x.mag, y.mag) });
}
\end{Verbatim}
\end{singlespace}
\end{minipage}\par\medskip

\par\noindent\begin{minipage}{\linewidth}
\fn{divS}{\code{n}, \code{m} --- делимое и делитель; \code{q}, \code{r} --- выходные параметры.}{\code{false} при нулевом делителе; иначе \code{true}, частное и неотрицательный остаток. Для отрицательного делимого с ненулевым остатком выполняется поправка евклидова деления.}
\begin{singlespace}
\captionof{listing}{\code{divS}}
\begin{Verbatim}
static bool divS(const SNum& n, const SNum& m, SNum& q, SNum& r) {
    if (isZero(m.mag)) return false;
    QR qr = divMag(n.mag, m.mag);
    Num Q = qr.q;
    Num rem = qr.r;
    if (n.neg && !isZero(rem)) {               // делимое отрицательное и не делится нацело:
        Q   = addMag(Q, std::string(1, ONE));  //   округляем частное «вниз»,
        rem = subMag(m.mag, rem);              //   остаток = |m| - остаток
    }
    q = fixZero(SNum{ n.neg != m.neg, Q });
    r = SNum{ false, rem };
    return true;
}
\end{Verbatim}
\end{singlespace}
\end{minipage}\par\medskip

\subsection{Ввод, вывод и обработка ошибок}

Выражение вводится в виде трёх разделённых пробелами частей: \code{<число> <операция> <число>}. Разрешены операции \code{+}, \code{-}, \code{*}, \code{/}. Знак операнда записывается слитно с числом, например \code{-bc / e}. Разбор выделяет ровно три части: дополнительные слова считаются ошибкой.

\par\noindent\begin{minipage}{\linewidth}
\fn{canon}{\code{ch} --- буква операнда после приведения к нижнему регистру.}{Символ $c$ для входного $d$; остальные символы сохраняются. Канонизация обеспечивает равносильность двух обозначений одной цифры.}
\begin{singlespace}
\captionof{listing}{\code{canon}}
\begin{Verbatim}
static char canon(char ch) { return ch == 'd' ? 'c' : ch; }
\end{Verbatim}
\end{singlespace}
\end{minipage}\par\medskip

\par\noindent\begin{minipage}{\linewidth}
\fn{parseOperand}{\code{tok} --- строка операнда; \code{out} --- выходное число.}{Пустая строка при успехе либо сообщение об ошибке. Обрабатываются необязательный знак, буквы $a,\ldots,h$ в обоих регистрах и отдельная запись \code{1} со знаком или без него, равная цифре $b$. Смешивать \code{1} с буквами нельзя.}
\begin{singlespace}
\captionof{listing}{\code{parseOperand}}
\begin{Verbatim}
static std::string parseOperand(const std::string& tok, SNum& out) {
    auto it = tok.begin();
    bool neg = false;
    if (it != tok.end() && (*it == '+' || *it == '-')) { neg = (*it == '-'); ++it; }
    std::string body(it, tok.end());
    if (body.empty()) return "после знака нет числа";
    if (body == "1") {                          // допустимая цифра: «+1» / «-1» (то есть b)
        out = { neg, std::string(1, ONE) };
        return "";
    }
    static const std::string LETTERS = "abcdefgh";
    Num mag;
    for (char raw : body) {
        char ch = (char)std::tolower((unsigned char)raw);
        if (std::isdigit((unsigned char)ch))
            return std::string("цифра '") + raw + "' недопустима (из цифр можно только 1 как +1 / -1)";
        if (LETTERS.find(ch) == std::string::npos)
            return std::string("символ '") + raw + "' недопустим (нужны буквы a b c d e f g h)";
        mag.push_back(canon(ch));
    }
    if (mag.size() > MAX_DIGITS) return "в числе больше 8 разрядов";
    mag = strip(mag);
    out = { neg && !isZero(mag), mag };
    return "";
}
\end{Verbatim}
\end{singlespace}
\end{minipage}\par\medskip

Функция проверяет наличие числа после знака, допустимость каждого символа и длину исходной буквенной записи. Затем заменяет $d$ на $c$, удаляет ведущие нули и нормализует знак. Например, \code{aabd} преобразуется в \code{bc}, а \code{-a} --- в \code{a}.

\par\noindent\begin{minipage}{\linewidth}
\fn{render}{\code{v} --- число со знаком; \code{cd} --- символ $c$ или $d$.}{Строковая запись числа с указанным обозначением максимальной цифры. В начало отрицательного числа добавляется знак минус.}
\begin{singlespace}
\captionof{listing}{\code{render}}
\begin{Verbatim}
static std::string render(const SNum& v, char cd) {
    std::string s = v.neg ? "-" : "";
    for (char ch : v.mag) s.push_back(ch == 'c' ? cd : ch);
    return s;
}
\end{Verbatim}
\end{singlespace}
\end{minipage}\par\medskip

\par\noindent\begin{minipage}{\linewidth}
\fn{both}{\code{v} --- число со знаком.}{Если модуль содержит $c$, выводятся две равные записи: со всеми $c$ и со всеми $d$. Например, \code{cc = dd}. Смешанные варианты \code{cd} и \code{dc} отдельно не перечисляются. Если $c$ нет, выводится одна запись.}
\begin{singlespace}
\captionof{listing}{\code{both}}
\begin{Verbatim}
static std::string both(const SNum& v) {
    if (v.mag.find('c') == std::string::npos) return render(v, 'c');
    return render(v, 'c') + " = " + render(v, 'd');
}
\end{Verbatim}
\end{singlespace}
\end{minipage}\par\medskip

\par\noindent\begin{minipage}{\linewidth}
\fn{warnOverflow}{\code{v} --- вычисленный результат.}{При длине модуля больше восьми выведено предупреждение о переполнении. Сам результат не обрезается и не заменяется остатком по модулю $7^8$.}
\begin{singlespace}
\captionof{listing}{\code{warnOverflow}}
\begin{Verbatim}
static void warnOverflow(const SNum& v) {
    if (v.mag.size() > MAX_DIGITS) std::cout << "   (внимание: результат длиннее 8 разрядов - переполнение разрядной сетки)\n";
}
\end{Verbatim}
\end{singlespace}
\end{minipage}\par\medskip

\par\noindent\begin{minipage}{\linewidth}
\fn{evalLine}{\code{line} --- строка выражения.}{Вычислен и выведен результат либо возвращено сообщение об ошибке формата или операндов. При делении выводятся частное и остаток. Сообщение о делении на ноль печатается непосредственно внутри функции.}
\begin{singlespace}
\captionof{listing}{\code{evalLine}}
\begin{Verbatim}
static std::string evalLine(const std::string& line) {
    std::istringstream in(line);
    std::string t1, op, t2, extra;
    if (!(in >> t1 >> op >> t2) || (in >> extra))
        return "формат: <число> <операция> <число>, между числами и операцией обязательны пробелы";
    if (op != "+" && op != "-" && op != "*" && op != "/")
        return "операция должна быть одной из: + - * /";

    SNum x, y;
    std::string err = parseOperand(t1, x);
    if (!err.empty()) return "первое число: " + err;
    err = parseOperand(t2, y);
    if (!err.empty()) return "второе число: " + err;

    if (op == "/") {
        SNum q, r;
        if (!divS(x, y, q, r)) {
            std::cout << "Ошибка: деление на a (ноль) невозможно.\n";
            return "";
        }
        std::cout << "Частное: " << both(q) << "\n";
        std::cout << "Остаток: " << both(r) << "\n";
        warnOverflow(q);
        return "";
    }
    SNum res = (op == "+") ? addS(x, y) : (op == "-") ? subS(x, y) : mulS(x, y);
    std::cout << "Результат: " << both(res) << "\n";
    warnOverflow(res);
    return "";
}
\end{Verbatim}
\end{singlespace}
\end{minipage}\par\medskip

\Needspace{20\baselineskip}
\subsection{Вывод таблиц}

\par\noindent\begin{minipage}{\linewidth}
\fn{printTable}{\code{title} --- заголовок; \code{t} --- таблица цифр.}{В консоль выведена таблица с заголовками строк и столбцов. Порядок строк и столбцов определяется контейнером \code{std::map}.}
\begin{singlespace}
\captionof{listing}{\code{printTable}}
\begin{Verbatim}
static void printTable(const std::string& title, const Table& t) {
    std::cout << "\n" << title << "\n   |";
    for (const auto& kv : t.begin()->second) std::cout << ' ' << kv.first;
    std::cout << "\n";
    for (const auto& row : t) {
        std::cout << " " << row.first << " |";
        for (const auto& cell : row.second) std::cout << ' ' << cell.second;
        std::cout << "\n";
    }
}
\end{Verbatim}
\end{singlespace}
\end{minipage}\par\medskip

\par\noindent\begin{minipage}{\linewidth}
\fn{printDiv}{Таблица \code{DIV}.}{В консоль выведены пары частного и остатка. Столбца для нулевого делителя нет.}
\begin{singlespace}
\captionof{listing}{\code{printDiv}}
\begin{Verbatim}
static void printDiv() {
    std::cout << "\nДеление (частное,остаток): строка - делимое, столбец - делитель\n   |";
    for (const auto& kv : DIV.begin()->second) std::cout << "  " << kv.first << " ";
    std::cout << "\n";
    for (const auto& row : DIV) {
        std::cout << " " << row.first << " |";
        for (const auto& cell : row.second) std::cout << ' ' << cell.second.first << ',' << cell.second.second;
        std::cout << "\n";
    }
}
\end{Verbatim}
\end{singlespace}
\end{minipage}\par\medskip

\par\noindent\begin{minipage}{\linewidth}
\fn{showTables}{Таблицы одноразрядных операций.}{Последовательно выведены сложение, вычитание, младшая и старшая цифры произведения, деление. Дополнительно сообщается о равносильности $d$ и $c$.}
\begin{singlespace}
\captionof{listing}{\code{showTables}}
\begin{Verbatim}
static void showTables() {
    printTable("Сложение (строка + столбец)", ADD);
    printTable("Вычитание (строка - столбец)", SUB);
    printTable("Умножение, младшая цифра (строка * столбец)", MUL);
    printTable("Умножение, старшая цифра (перенос)", MULHI);
    printDiv();
    std::cout << "\nЭлемент d равносилен c.\n";
}
\end{Verbatim}
\end{singlespace}
\end{minipage}\par\medskip

\Needspace{20\baselineskip}
\subsection{Меню и функция main}

\par\noindent\begin{minipage}{\linewidth}
\fn{calcLoop}{Поток пользовательского ввода.}{Выполнена последовательность вычислений. После ошибки выводится пояснение, и ввод повторяется. Команда \code{m} или \code{M} возвращает в главное меню; конец потока завершает цикл.}
\begin{singlespace}
\captionof{listing}{\code{calcLoop}}
\begin{Verbatim}
static void calcLoop() {
    std::string line;
    while (true) {
        std::cout << "\nВведите выражение (например: b + -a, -a + b, c * 1, bc / e), m - в меню: ";
        if (!std::getline(std::cin, line)) return;
        if (line == "m" || line == "M") return;
        std::string err = evalLine(line);
        if (!err.empty()) std::cout << "Ошибка ввода: " << err << ". Введите снова.\n";
    }
}
\end{Verbatim}
\end{singlespace}
\end{minipage}\par\medskip

\par\noindent\begin{minipage}{\linewidth}
\fn{main}{Стандартные потоки ввода и вывода.}{Сначала построен порядок цифр, затем организовано главное меню: \code{1} --- вычисление выражений, \code{2} --- таблицы, \code{0} --- выход. Неизвестный пункт вызывает сообщение с допустимыми вариантами.}
\begin{singlespace}
\captionof{listing}{\code{main}}
\begin{Verbatim}
int main() {
    buildOrder();
    std::string line;
    while (true) {
        std::cout << "\n=== Конечная арифметика, вариант X8 ===\n"
                     "1 - вычислить выражение\n"
                     "2 - показать таблицы\n"
                     "0 - выход\n"
                     "Выбор: ";
        if (!std::getline(std::cin, line)) break;
        if (line == "1")      calcLoop();
        else if (line == "2") showTables();
        else if (line == "0") break;
        else std::cout << "Неизвестный пункт меню, введите 1, 2 или 0.\n";
    }
    return 0;
}
\end{Verbatim}
\end{singlespace}
\end{minipage}\par\medskip

\Needspace{20\baselineskip}
\subsection{Оценка вычислительной сложности}

Размер алфавита постоянен и равен семи, поэтому обращения к таблицам и сравнение отдельных цифр рассматриваются как операции постоянной стоимости. Пусть $n,m$ --- длины модулей операндов, $k=\max(n,m)$.

Сравнение модулей требует $O(k)$ действий. В сложении и вычитании обрабатывается $O(k)$ разрядов, однако каждая вставка в начало \code{std::string} может перемещать уже накопленные символы. Поэтому для данной реализации верхняя оценка времени равна $O(k^2)$.

При умножении выполняется до $m$ частичных произведений и сложений длины не более $n+m$. Для операндов сравнимой длины консервативная оценка составляет $O(k^3)$. Деление обрабатывает $n$ цифр делимого; для каждой цифры частного выполняется не более шести вычитаний. С учётом вставок в начало строк его верхняя оценка при сравнимых длинах также составляет $O(k^3)$.

Рабочая память для строковых операндов и результатов составляет $O(n+m)$; таблицы имеют постоянный размер. При входном ограничении в восемь разрядов эти затраты малы. Возможное улучшение --- накапливать результат от младших разрядов через \code{push_back}, а затем переворачивать строку; это устранило бы повторные вставки в её начало.

\section{Результаты работы}

\subsection{Контрольные примеры вычислений}

Программа собрана компилятором \code{g++} с параметрами \code{-std=c++17 -O2 -Wall -Wextra}. Компиляция завершилась без предупреждений. В таблице~\ref{tab:examples} приведены результаты фактических запусков. Десятичные значения используются только для независимой проверки и не выводятся самой программой.

\begin{table}[H]
\caption{Проверка основных арифметических операций}\label{tab:examples}
\centering\small
\begin{tabularx}{\textwidth}{|l|L|L|}\hline
Выражение & Результат программы & Проверка\\\hline
\code{c + b} & \code{ba} & $6+1=7$\\\hline
\code{ba - b} & \code{c = d} & $7-1=6$\\\hline
\code{c * c} & \code{gb} & $6\cdot6=36$\\\hline
\code{e - c} & \code{-f} & $2-6=-4$\\\hline
\code{-e + g} & \code{h} & $-2+5=3$\\\hline
\code{bc / e} & \code{q: c = d}, \code{r: b} & $13=6\cdot2+1$\\\hline
\code{-bc / e} & \code{q: -ba}, \code{r: b} & $-13=(-7)\cdot2+1$\\\hline
\code{bc / -e} & \code{q: -c = -d}, \code{r: b} & $13=(-6)\cdot(-2)+1$\\\hline
\code{-bc / -e} & \code{q: ba}, \code{r: b} & $-13=7\cdot(-2)+1$\\\hline
\end{tabularx}
\end{table}

Ниже приведены сокращённые консольные протоколы. Введённые выражения помечены символом \code{>}; повторяющиеся приглашения программы опущены. Текст результатов соответствует выводу исполняемого файла.

\begin{figure}[H]
\begin{singlespace}
\begin{Verbatim}[numbers=none,fontsize=\small]
=== Конечная арифметика, вариант X8 ===
1 - вычислить выражение
2 - показать таблицы
0 - выход

> c + b
Результат: ba
> ba - b
Результат: c = d
> c * c
Результат: gb
> -bc / e
Частное: -ba
Остаток: b
\end{Verbatim}
\end{singlespace}
\caption{Основные операции и деление отрицательного числа}
\end{figure}

\subsection{Нормализация и предельные значения}

Ввод \code{D + b} дал результат \code{ba}: заглавная буква приведена к нижнему регистру, после чего $d$ заменено на $c$. Для \code{bd + b} получено \code{ea}, что соответствует равенству $13+1=14$. В выражении \code{-a + a} результат равен \code{a}, отрицательный ноль не появляется.

При сложении максимального восьмиразрядного числа с единицей получен девятиразрядный результат \code{baaaaaaaa}, соответствующий $7^8=5\,764\,801$. Программа сохранила полный результат и вывела предупреждение. При умножении двух максимальных операндов получена шестнадцатиразрядная запись:
\[
V(cccccccgaaaaaaab)=(7^8-1)^2.
\]

\begin{figure}[H]
\begin{singlespace}
\begin{Verbatim}[numbers=none,fontsize=\small]
> D + b
Результат: ba
> -a + a
Результат: a
> cccccccc + b
Результат: baaaaaaaa
   (внимание: результат длиннее 8 разрядов -
   переполнение разрядной сетки)
> cccccccc * cccccccc
Результат: cccccccgaaaaaaab = dddddddgaaaaaaab
   (внимание: результат длиннее 8 разрядов -
   переполнение разрядной сетки)
\end{Verbatim}
\end{singlespace}
\caption{Нормализация и переполнение разрядной сетки}
\end{figure}

\Needspace{20\baselineskip}
\subsection{Обработка некорректного ввода}

Проверены деление на ноль, неизвестная операция, отсутствие пробелов, недопустимая цифра и слишком длинный операнд. Во всех случаях программа сообщила причину ошибки и сохранила возможность следующего ввода.

\begin{table}[H]
\caption{Проверка ошибочного ввода}
\centering\small
\begin{tabularx}{\textwidth}{|l|L|}\hline
Ввод & Реакция программы\\\hline
\code{b / a} & Сообщение о невозможности деления на $a$ (ноль)\\\hline
\code{aaaaaaaaa + b} & В первом числе больше восьми разрядов\\\hline
\code{b+b} & Сообщение о формате и обязательных пробелах\\\hline
\texttt{b \% e} & Сообщение о допустимых операциях \code{+ - * /}\\\hline
\code{b2 + a} & Цифра \code{2} недопустима\\\hline
\code{i + a} & Буква \code{i} не входит в алфавит\\\hline
\end{tabularx}
\end{table}

\subsection{Автоматическая проверка}

Для независимой проверки подготовлена программа \code{course-work/build/verify_report.cpp}. Она подключает функции исходного файла и сравнивает результаты с арифметикой стандартного целочисленного типа. Перевод чисел используется только при проверке.

Проверены все 49 пар цифр для каждой из таблиц \code{ADD}, \code{SUB}, \code{MUL}, \code{MULHI} и 42 пары с ненулевым делителем для \code{DIV}: всего 238 элементов.

Перебраны все пары чисел от $-48$ до $48$ включительно (9409 пар) и 10\,000 псевдослучайных пар из диапазона $[-5\,764\,800;5\,764\,800]$. Генератор \code{std::mt19937_64} инициализирован значением 20261116.

Для каждой из 19\,409 пар проверены четыре операции: всего 77\,636 проверок. Для деления проверялись равенство $n=qm+r$, границы $0\leq r<|m|$ либо отказ при нулевом делителе. Дополнительно проверены 14 вариантов разбора операндов, перенос в девятый разряд и произведение максимальных модулей. Все проверки завершились успешно.

\structsection{ЗАКЛЮЧЕНИЕ}

В результате курсовой работы разработана программа конечной арифметики для варианта X8. По заданному правилу перехода определены семь различных значений цифр, установлен их порядок и составлены таблицы сложения, вычитания, умножения и деления. Символы $c$ и $d$ обрабатываются как два обозначения одной цифры.

Реализованы многоразрядные операции над числами со знаком. Сложение учитывает переносы, вычитание --- заимствования, умножение выполняется по частичным произведениям, деление возвращает частное и неотрицательный остаток. Вычисления выполняются над строковыми записями с использованием таблиц цифр.

Программа предоставляет меню, просмотр таблиц и повторный ввод выражений. Обрабатываются недопустимые символы, ошибки формата, превышение длины входного числа и деление на ноль. Если результат содержит более восьми разрядов, выводится предупреждение и сохраняется полная запись результата.

Корректность проверена на всех элементах таблиц, контрольных примерах, полном переборе малых операндов и псевдослучайных парах в допустимом диапазоне. Все выполненные проверки прошли успешно.

Недостатком реализации являются повторные вставки в начало строки при построении результатов. При увеличении допустимой разрядности это приводит к дополнительным перемещениям символов. Дальнейшее развитие программы возможно за счёт накопления разрядов в обратном порядке, выделения арифметического ядра в отдельный модуль и добавления других операций, например возведения в степень.

\structsection{СПИСОК ИСПОЛЬЗОВАННЫХ ИСТОЧНИКОВ}

\begin{enumerate}
\item Дужнов, Р. Ю. Конечная арифметика, вариант X8 : исходный код программы на языке C++. --- 2026. --- Файл \code{course-work/main.cpp}.
\item Дужнов, Р. Ю. Проверка конечной арифметики, вариант X8 : вспомогательная программа. --- 2026. --- Файл \code{course-work/build/verify_report.cpp}.
\end{enumerate}

\end{document}
